% AUTORIDAD DE FUENTE: Exposé IX §13.3, líneas francesas 4035--4173;
% índices de escaneo 500--504, folios 489--493.

\medskip
\noindent
13.3. \underline{La extensión de Raynaud analítica}.

Estamos ya en condiciones de parafrasear, en geometría analítica, la
construcción de la extensión de Raynaud dada en 7. Supongamos que $A$ es
como en 13.1, con $A=\sisi{A_o}{A^o}$, y consideremos $R_o=R_s$, que es
un $\mathbb Z$-módulo libre de tipo finito. Tras restringir $S$ a un
entorno abierto suficientemente pequeño de $s$ e interpretar $R_o$ como
el grupo de gérmenes de secciones de $R$ cerca de $s$, se obtiene un
homomorfismo

\begin{equation}\tag{13.3.1}
 R_{oS}\longrightarrow R\subset E,
\end{equation}

\noindent
donde $R_{oS}$ denota el grupo analítico constante sobre $S$ con valor
$R_o$. Eligiendo el entorno de $s$ lo bastante pequeño para que toda
sección de $\underline{\mathcal O}_S$ cuyo germen en $s$ sea cero sea
ella misma cero, se ve que el homomorfismo precedente es inyectivo. Su
imagen es, por tanto, un subgrupo analítico constante de $R$; es
manifiestamente el mayor subgrupo constante de $R$, se llama la
\underline{parte fija} de $R$ y se denota por $R^f$:

%%%%%%%%%%  scan idx 501  |  folio 490  %%%%%%%%%%

\begin{equation}\tag{13.3.2}
 R_{oS}\simeq R^f\hookrightarrow R,
 \qquad\text{e induce un isomorfismo }R_s^f\simeq R_s.
\end{equation}

\noindent
Podemos considerar entonces el grupo cociente

\begin{equation}\tag{13.3.3}
 A^\natural=E/R^f,
\end{equation}

\noindent
y obtener, teniendo en cuenta (13.1.3), una sucesión exacta canónica

\begin{equation}\tag{13.3.4}
 0\longrightarrow R/R^f\stackrel{i}{\longrightarrow}A^\natural
 \longrightarrow A\longrightarrow0.
\end{equation}

\noindent
Obsérvese que, tras restringir $S$ a un entorno de $s$, $R^f$ es siempre
cerrado en $E$, de modo que $A^\natural$ es separado sobre $S$.
Obsérvese asimismo que el homomorfismo
$A^\natural\longrightarrow A$ es étale y, por tanto, induce un
isomorfismo

\begin{equation}\tag{13.3.5}
 \underline{\operatorname{Lie}}(A^\natural)
 \stackrel{\sim}{\longrightarrow}
 \underline{\operatorname{Lie}}(A)=\underline E.
\end{equation}

\noindent
Además, para cada $n\geqslant0$, el homomorfismo canónico
$A^\natural\longrightarrow A$ induce un isomorfismo
$A_n^\natural\simeq A_n$ (pues
$R_n^f\stackrel{\sim}{\longrightarrow}R_n$) y, en consecuencia, un
isomorfismo

\begin{equation}\tag{13.3.6}
 \widehat{A^\natural}\stackrel{\sim}{\longrightarrow}\widehat A.
\end{equation}

\noindent
Este último hecho muestra que $A^\natural$ desempeña efectivamente el
papel del grupo, denotado también por $A^\natural$, que consideramos en
7. Aquí se ha construido sin ninguna hipótesis sobre la fibra
$A_s=A_o$ de la forma «extensión de un grupo abeloide por un toro». Tal
hipótesis será, no obstante, necesaria si queremos recuperar sobre
$A^\natural$ la importante estructura de extensión (7.2.3). En cambio,
el homomorfismo $A^\natural\longrightarrow A$ de (13.3.4) es de
naturaleza esencialmente trascendente y no tiene análogo en la teoría
puramente algebraica de §~7. En §~14 volveremos a la cuestión de una
interpretación puramente algebraica del homomorfismo
$i:R/R^f\longrightarrow A^\natural$ de (13.3.4).

%%%%%%%%%%  scan idx 502  |  folio 491  %%%%%%%%%%

Supongamos primero que se nos da un subtoro complejo

\begin{equation}\tag{13.3.7}
 T_o\subset A_o,
\end{equation}

\noindent
donde por «\underline{toro complejo}» entendemos un grupo analítico
isomorfo a una potencia cartesiana de $\mathbb C^*$. Sea

\begin{equation}\tag{13.3.8}
 T=T_o\times S
\end{equation}

\noindent
el toro constante sobre $S$ con valor $T_o$. Procediendo como en 5.1,
se obtiene un único homomorfismo (de hecho, una inmersión) de sistemas
inductivos de grupos analíticos

\begin{equation*}
 \widehat T\longrightarrow\widehat{A^\natural},
\end{equation*}

\noindent
que sobre $S_o$ se reduce a (13.3.7). Por 13.2.2, este homomorfismo
procede de un único germen de homomorfismos

\begin{equation}\tag{13.3.9}
 T\longrightarrow A^\natural,
\end{equation}

\noindent
y, tomando $S$ suficientemente pequeño alrededor de $s$, podemos
suponer que está definido sobre todo $S$. El homomorfismo correspondiente

\begin{equation*}
 \underline{\operatorname{Lie}}(T)\longrightarrow
 \underline{\operatorname{Lie}}(A^\natural)\simeq\underline E
\end{equation*}

\noindent
es inyectivo sobre cada $S_n$, luego sobre $\widehat S$ y, por tanto,
sobre $\underline{\mathcal O}_s$ por descenso plano. Existe, pues, un
entorno de $s$ en $S$, que podemos tomar igual a todo $S$, sobre el cual
$\underline{\operatorname{Lie}}(T)$ se identifica con un sumando
localmente directo $\underline E^t$ de $\underline E$; denotemos por
$E^t$ el subfibrado vectorial correspondiente de $E$:

%%%%%%%%%%  scan idx 503  |  folio 492  %%%%%%%%%%

\begin{equation}\tag{13.3.10}
 \underline{\operatorname{Lie}}(T)\xrightarrow{\ \sim\ }
 \underline E^t\subset\underline E\simeq
 \underline{\operatorname{Lie}}(A^\natural)\simeq
 \underline{\operatorname{Lie}}(A),
 \qquad E^t\subset E.
\end{equation}

\noindent
Asimismo, sean

\begin{equation}\tag{13.3.11}
 R_o^t\subset R_o,
 \qquad R^t=R_{oS}^t\subset R_{oS}\simeq R^f
\end{equation}

\noindent
respectivamente la subred de $R_o$ definida por el subtoro $T_o$ de
$A_o$ y el grupo analítico constante correspondiente sobre $S$. Entonces
el homomorfismo (13.3.9) corresponde simplemente al homomorfismo

\begin{equation*}
 E^t/R_{oS}^t=E^t/R^t\longrightarrow E/R_{oS}=E/R^f
\end{equation*}

\noindent
inducido por las inclusiones

\begin{equation*}
 E^t\hookrightarrow E,
 \qquad R^t\hookrightarrow R^f.
\end{equation*}

\noindent
Utilizando el hecho de que (13.3.9) induce una inmersión en la fibra
sobre $s$, concluimos que el homomorfismo

\begin{equation*}
 R_o/R_o^t\longrightarrow E_o/E_o^t=(E/E^t)_o
\end{equation*}

\noindent
es inyectivo. Se sigue inmediatamente que

\begin{equation*}
 R_o/R_o^t\longrightarrow(E/E^t)_{s'}
\end{equation*}

\noindent
sigue siendo inyectivo para $s'$ próximo a $s$. Así, tras restringir $S$
alrededor de $s$, podemos suponer que (13.3.9) \underline{es una
inmersión cerrada}.

%%%%%%%%%%  scan idx 504  |  folio 493  %%%%%%%%%%

Nos vemos, pues, llevados a introducir

\begin{equation}\tag{13.3.12}
 B=A^\natural/T\simeq(E/E^t)/(R^f/R^t),
\end{equation}

\noindent
de modo que $A^\natural$ adquiere una estructura de extensión

\begin{equation}\tag{13.3.13}
 0\longrightarrow T\longrightarrow A^\natural
 \longrightarrow B\longrightarrow0,
\end{equation}

\noindent
donde $B$ es un grupo analítico relativo sobre $S$, liso y separado
sobre $S$, de fibras conexas. Por construcción tenemos, pues,

\begin{equation}\tag{13.3.14}
 \underline{\operatorname{Lie}}(B)\simeq
 \underline E/\underline E^t,
\end{equation}

\noindent
y la sucesión exacta exponencial de $B$ es

\begin{equation}\tag{13.3.15}
 0\longrightarrow R^f/R^t\longrightarrow E/E^t
 \longrightarrow B\longrightarrow0.
\end{equation}

\noindent
Obsérvese que

\begin{equation}\tag{13.3.16}
 R^f/R^t=(R_o/R_o^t)_S
\end{equation}

\noindent
es un grupo analítico constante sobre $S$. Se sigue de 13.1.5 que $B$
es un grupo abeloide sobre $S$ en un entorno de $s$ si y solo si su fibra
en $s$,

\begin{equation*}
 B_o=A_o/T_o,
\end{equation*}

\noindent
es un grupo analítico abeloide; en ese caso, $A_o$ aparece como una
extensión de tal grupo $B_o$ por un toro $T_o$. \underline{Toda manera de
escribir $A_o$ como una extensión de un grupo abeloide por un toro
complejo} $(*)$ \underline{da lugar, por tanto, sobre $A$, en un entorno de} $s$, \underline{a
una única estructura de extensión} (13.3.13) \underline{de un grupo
abeloide $B$ por un toro $T$, que induce la estructura de extensión dada
en $s$}.
